\documentclass[../main.tex]{subfiles}

\graphicspath{{\subfix{../images/}}}


\begin{document}

\newpage


\section{General introduction}


This section introduces the course organization, illustrating its structure, how it will be structured, the exam rules, and the main aspects that characterize it as a whole.

\subsection{Course structure}

The course is divided into two parts, each taught by a different instructor.
The first part covers topics that are more hardware-oriented, specifically:

\begin{enumerate}
	\item[\textbf{-}] Measurements and Instruments – 1 CFU
	\item[\textbf{-}] Operational Amplifiers – 2 CFU
	\item[\textbf{-}] Power Supplies – 1 CFU
\end{enumerate}

The second part of the course focuses on more application-oriented and integration aspects, including:

\begin{enumerate}
	\item[\textbf{-}] Firmware – 2 CFU
	\item[\textbf{-}] Sensors - 1 CFU
	\item[\textbf{-}] Actuators and Motors – 1 CFU
\end{enumerate}

There is no clear separation between theoretical lessons and practical exercises: both activities take place within the same timetable and are integrated with each other. There are also specific laboratory slots, during which practical exercises will be developed, according to the calendar communicated in class.

\subsection{Material}

As this is the first edition of the course, no preliminary material has been made available. Furthermore, the lessons will not be videotaped and only the transparencies containing what the teacher writes and explains during the classroom sessions will be provided.\\

Despite these limitations, for the purposes of this paper, we have chosen to refer to a number of texts that we consider useful for exploring and presenting the main concepts covered in the course in a clear and structured manner. These are the two textbooks used:

\begin{enumerate}
	\item[\textbf{-}] Sedra, Smith, Microelectronic Circuits, 7th edition, Oxford University Press, 2016.
	\item[\textbf{-}] Jaeger, Blalock, Microelectronic circuit design, 6th edition, McGraw-Hill, 2022
\end{enumerate}

In addition to the books reported before material such as user guide, datasheets, reference manuals, application notes and more general technical literature will be provided when needed.

\subsection{Lessons exercises and laboratories}

As said before there is no clear division between lessons and exercises sessions. The only well defined activity will be each of the six laboratories covering the topic discussed in the previous weeks.
The attendance to laboratories is not mandatory but strongly suggested. Each laboratory session requires a mandatory report to be submitted in a one week time from the laboratory session. The report will be submitted as a group and will be evaluated with a maximum of 2 points each for a total of 12.

\subsection{Exam}

The course grade is determined by the sum of the scores obtained in the laboratory reports, for a maximum of 12 points, and the score obtained in the compulsory oral exam, which is worth up to 18 points.
The oral exam is conducted individually and lasts approximately 45 minutes. During the interview, theoretical knowledge, the ability to solve short exercises, and comprehension of the laboratory reports submitted during the semester will be assessed.\\

Passing this exam is closely linked to the final exam of the Computer Engineering Degree Course at the Polytechnic University of Turin: passing the course also implies passing the relevant final exam.

The final exam consists of writing an additional laboratory report, which will require the ability to develop a basic project by applying the knowledge acquired during the course.\\

As this is the first edition of the course, the information provided in this section may be subject to change or may not be fully up to date. It is therefore recommended that you always refer to the instructions provided by the instructor and to the information published on the official course portal.









\end{document}